\FloatBarrier
\subsection{Four configurations and what their contrasts isolate}
\label{sec:r42}
Filter switch $ f $ and code switch $ c $ vary while sampling, Gaussian/BSC
models, quantiser, framing and receiver remain common. Here $ x $ contains
32000 transmission samples; $\mathcal R $ returns 96000 reference-grid
positions and $ d_{fc}$ denotes their full-record NMSE:
\begin{align}
 y&=x+\eta,\quad z_f=F_f(y),\quad
 b_f=\operatorname{pack}(\mathcal Q(z_f)),\nonumber\\
 t_{fc}&=E_c(b_f),\quad r_{fc}=t_{fc}\oplus\varepsilon_c,\nonumber\\
 \widehat x_{fc}&=\mathcal R\!\left[
 \mathcal Q^{-1}\!\left(\operatorname{unpack}(D_c(r_{fc}))\right)\right].
 \label{eq:r42-chain}
\end{align}
Here $ F_0=E_0=D_0=I $; $ F_1 $ is the frozen FIR and $ E_1,D_1 $ the taught
Hamming(7,4). Neither bypasses the protective modules, not the degradations.
Code-only quantises noisy unfiltered samples; filter-only sends unprotected
bits through the BSC. The receiver has only received bits and public
codebook/framing/output-length metadata, never the reference or true masks.

\begin{table}[htbp]
 \centering
 \small
 \caption{Four actual paths carrying 448000 real information bits. Gaussian noise, quantisation, BSC and reconstruction remain enabled in every configuration.}
 \label{tab:r42-configurations}
 \begin{tabular}{llrrr}
\toprule
ID & Configuration & $f$ & $c$ & Sent bits / logical kbit/s \\
\midrule
00 & neither & 0 & 0 & 448000 / 224 \\
10 & filter-only & 1 & 0 & 448000 / 224 \\
01 & code-only & 0 & 1 & 784000 / 392 \\
11 & both & 1 & 1 & 784000 / 392 \\
\bottomrule
\end{tabular}

 
 
\end{table}

One Gaussian vector is shared within each $(S,r)$ across all four paths
and $ p $ values. Same-code filter pairs replay the same BSC leaf and shape;
uncoded/coded layouts have independent masks and unequal lengths.
\begin{align}
 \text{Gaussian}:&(1,4,a,\mathrm{SNR\ ID},r),\nonumber\\
 \text{BSC}:&(2,4,a,\mathrm{joint\ ID},r,\mathrm{layout\ ID}).
 \label{eq:r42-rng}
\end{align}
At $ p=0 $ no mask is drawn. The actual decoder's translation identity
(Eq.~\eqref{eq:r42-pair-invariant}) gives equal XOR error patterns/BER
within 00/10 and 01/11, while their waveform NMSE can differ through
bit values, upstream errors and amplitude-error signs.
The conditional reductions and interaction are
\begin{align}
 g_{F\mid c=0}&=d_{00}-d_{10},&
 g_{F\mid c=1}&=d_{01}-d_{11},\nonumber\\
 g_{C\mid f=0}&=d_{00}-d_{01},&
 g_{C\mid f=1}&=d_{10}-d_{11},\nonumber\\
 I&=d_{00}-d_{10}-d_{01}+d_{11}.
 \label{eq:r42-effects}
\end{align}
Positive reductions improve quality; negative values are retained.
Comparisons at fixed $ c $ isolate enabling the filter; those at fixed $ f $
isolate enabling ECC including its extra transmitted bits. Comparing only
00 with 11 confounds these effects. Interaction is not a responsibility
percentage, and these switches do not separately identify changing the
sampling rate or bit depth. Rates remain 224/392 kbit/s, not equal energy.

Twenty new pilot bundles over 22 unique conditions produce 1760 paths.
All are independently replayed; mask invariants and clean/identity controls
pass. Pilot variance alone plans 1400 validation bundles for 88 condition
means plus four main-objective means, targeting 5 \% relative half-width.
The smallest passing multiple of 50 is frozen before phase 1
(Eq.~\eqref{eq:r42-planning}); a pilot prediction is not a guarantee of
achieved precision. No pilot observation enters the final means.
